% GENERATED FILE. Edit canonical lessons, then run npm run generate:books.
% canonical-source-hash: fnv1a64:e218f05db0588114
% canonical-lessons: ZH-C160-saodi, ZH-C160-xiwan, ZH-C160-puchuang, ZH-C160-ganhuo, ZH-C160-zhaogu

\chapter{Doing the Chores}
\label{ch:zh-doing-the-chores}
\hlchaptermodality{160}

\begin{chapteropening}
\textbf{By the end of this chapter:} \emph{I can talk about sweeping, washing up, making the bed, doing jobs and looking after someone.}

Complete the last lesson of chapter 160: I can talk about sweeping, washing up, making the bed, doing jobs and looking after someone.
\end{chapteropening}

\section[sǎodì]{sǎodì — to sweep the floor}

\label{lesson:ZH-C160-saodi}



\begin{quote}
Before the new one: say the Mandarin for a balcony, then the Mandarin for a window.
\end{quote}

\subsection*{You'll want to know: sǎodì}
\emph{sǎodì} — \textquotedblleft{}to sweep the floor\textquotedblright{}.

\begin{quote}
\emph{wǒ sǎodì nǐ xǐ cài} — I'll sweep the floor; you wash the vegetables.
\end{quote}

Its characters come later; for now it is a word for the ear.

\begin{cousinweb}
\emph{sǎodì} is built from \zh{扫} 'to sweep' and \emph{dì} 'ground', both inherited from Middle Chinese, \emph{sawX} and \emph{dijH} (Guangyun, in Baxter's transcription).
\end{cousinweb}

\begin{grammarlens}[title={What you've built}]
One more word for saying what I do.
\end{grammarlens}

\subsection*{Your turn}
\begin{itemize}
\raggedright
  \item \emph{Say it:} \emph{sǎodì}
  \item \emph{Say it:} \emph{sǎodì}, once more
  \item \emph{Say it:} \emph{wǒ sǎodì nǐ xǐ cài} — I'll sweep the floor; you wash the vegetables.
  \item \emph{Recall:} say the Mandarin for a balcony, then the Mandarin for a window, then say \emph{sǎodì} again
\end{itemize}

\begin{tcolorbox}[breakable,colback=teal!4,colframe=teal!35!black,title={Before you move on}]
What does sǎodì mean? (\textquotedblleft{}To sweep the floor\textquotedblright{}.) Say it once more.
\end{tcolorbox}
\section[xǐwǎn]{xǐwǎn — to do the washing-up; to wash the dishes}

\label{lesson:ZH-C160-xiwan}



\begin{quote}
Before the new one: say the Mandarin for a window, then the Mandarin for to sweep the floor.
\end{quote}

\subsection*{You'll want to know: xǐwǎn}
\emph{xǐwǎn} — \textquotedblleft{}to do the washing-up; to wash the dishes\textquotedblright{}.

\begin{quote}
\emph{chī fàn yǐhòu wǒ xǐwǎn} — After we eat, I'll do the dishes.
\end{quote}

Its characters come later; for now it is a word for the ear.

Two third tones: xǐwǎn is said xíwǎn, the first rising.

\begin{cousinweb}
\emph{xǐwǎn} is built from \zh{洗} 'to wash' and \emph{wǎn} 'bowl', both inherited from Middle Chinese, \emph{sejX} and \emph{'wanX} (Guangyun, in Baxter's transcription), which writes the second another way and glosses it 'a vessel'.
\end{cousinweb}

\begin{grammarlens}[title={What you've built}]
One more word for saying what I do.
\end{grammarlens}

\subsection*{Your turn}
\begin{itemize}
\raggedright
  \item \emph{Say it:} \emph{xǐwǎn}
  \item \emph{Say it:} \emph{xǐwǎn}, once more
  \item \emph{Say it:} \emph{chī fàn yǐhòu wǒ xǐwǎn} — After we eat, I'll do the dishes.
  \item \emph{Recall:} say the Mandarin for a window, then the Mandarin for to sweep the floor, then say \emph{xǐwǎn} again
\end{itemize}

\begin{tcolorbox}[breakable,colback=teal!4,colframe=teal!35!black,title={Before you move on}]
What does xǐwǎn mean? (\textquotedblleft{}To do the washing-up; to wash the dishes\textquotedblright{}.) Say it once more.
\end{tcolorbox}
\section[pūchuáng]{pūchuáng — to make the bed}

\label{lesson:ZH-C160-puchuang}



\begin{quote}
Before the new one: say the Mandarin for to sweep the floor, then the Mandarin for to do the washing-up.
\end{quote}

\subsection*{You'll want to know: pūchuáng}
\emph{pūchuáng} — \textquotedblleft{}to make the bed\textquotedblright{}.

\begin{quote}
\emph{qǐchuáng yǐhòu bié wàng le pūchuáng} — Don't forget to make your bed when you get up.
\end{quote}

Its characters come later; for now it is a word for the ear.

\begin{cousinweb}
\emph{pūchuáng} is built from \emph{pū} 'to spread out' and \emph{chuáng} 'bed', both inherited from Middle Chinese, \emph{phu} and \emph{dzrjang} (Guangyun, in Baxter's transcription): to spread the bedding out.
\end{cousinweb}

\begin{grammarlens}[title={What you've built}]
One more word for saying what I do.
\end{grammarlens}

\subsection*{Your turn}
\begin{itemize}
\raggedright
  \item \emph{Say it:} \emph{pūchuáng}
  \item \emph{Say it:} \emph{pūchuáng}, once more
  \item \emph{Say it:} \emph{qǐchuáng yǐhòu bié wàng le pūchuáng} — Don't forget to make your bed when you get up.
  \item \emph{Recall:} say the Mandarin for to sweep the floor, then the Mandarin for to do the washing-up, then say \emph{pūchuáng} again
\end{itemize}

\begin{tcolorbox}[breakable,colback=teal!4,colframe=teal!35!black,title={Before you move on}]
What does pūchuáng mean? (\textquotedblleft{}To make the bed\textquotedblright{}.) Say it once more.
\end{tcolorbox}
\section[gànhuó]{gànhuó — to get the jobs done; to work}

\label{lesson:ZH-C160-ganhuo}



\begin{quote}
Before the new one: say the Mandarin for to do the washing-up, then the Mandarin for to make the bed.
\end{quote}

\subsection*{You'll want to know: gànhuó}
\emph{gànhuó} — \textquotedblleft{}to get the jobs done; to work\textquotedblright{}.

\begin{quote}
\emph{nǐmen dōu zài gànhuó ma} — Are you all busy working?
\end{quote}

Its characters come later; for now it is a word for the ear.

\begin{cousinweb}
\emph{gànhuó} is built from \emph{gàn} and \emph{huó}, both inherited from Middle Chinese, \emph{kanH} and \emph{hwat} (Guangyun, in Baxter's transcription). The dictionary glosses the first 'a stem' and 'strong' and the second 'alive'; how the pair came to mean doing the work is not settled.
\end{cousinweb}

\begin{grammarlens}[title={What you've built}]
One more word for saying what I do.
\end{grammarlens}

\subsection*{Your turn}
\begin{itemize}
\raggedright
  \item \emph{Say it:} \emph{gànhuó}
  \item \emph{Say it:} \emph{gànhuó}, once more
  \item \emph{Say it:} \emph{nǐmen dōu zài gànhuó ma} — Are you all busy working?
  \item \emph{Recall:} say the Mandarin for to do the washing-up, then the Mandarin for to make the bed, then say \emph{gànhuó} again
\end{itemize}

\begin{tcolorbox}[breakable,colback=teal!4,colframe=teal!35!black,title={Before you move on}]
What does gànhuó mean? (\textquotedblleft{}To get the jobs done; to work\textquotedblright{}.) Say it once more.
\end{tcolorbox}
\section[zhàogù]{zhàogù — to look after; to take care of}

\label{lesson:ZH-C160-zhaogu}



\begin{quote}
Before the new one: say the Mandarin for to make the bed, then the Mandarin for to get the jobs done.
\end{quote}

\subsection*{You'll want to know: zhàogù}
\emph{zhàogù} — \textquotedblleft{}to look after; to take care of\textquotedblright{}.

\begin{quote}
\emph{māma bìng le wǒ yào zhàogù tā} — Mum is ill, so I'm going to look after her.
\end{quote}

Its characters come later; for now it is a word for the ear.

\begin{cousinweb}
\emph{zhàogù} is built from \emph{zhào} 'bright, to shine' and \emph{gù} 'to turn and look'. The first is inherited from Middle Chinese \emph{tsyewH} (Guangyun, in Baxter's transcription); the second is listed there as \emph{kuH}, glossed 'to look back' and 'to care for'.
\end{cousinweb}

\begin{grammarlens}[title={What you've built}]
One more word for saying what I do.
\end{grammarlens}

\subsection*{Your turn}
\begin{itemize}
\raggedright
  \item \emph{Say it:} \emph{zhàogù}
  \item \emph{Say it:} \emph{zhàogù}, once more
  \item \emph{Say it:} \emph{māma bìng le wǒ yào zhàogù tā} — Mum is ill, so I'm going to look after her.
  \item \emph{Recall:} say the Mandarin for to make the bed, then the Mandarin for to get the jobs done, then say \emph{zhàogù} again
\end{itemize}

\begin{tcolorbox}[breakable,colback=teal!4,colframe=teal!35!black,title={Before you move on}]
What does zhàogù mean? (\textquotedblleft{}To look after; to take care of\textquotedblright{}.) Say it once more.
\end{tcolorbox}
